% Generated draft: requires mathematical and source review.
\documentclass[11pt]{article}
\usepackage[margin=25mm]{geometry}
\usepackage{fontspec}
\setmainfont{Latin Modern Roman}
\usepackage{amsmath,amssymb,mathtools,unicode-math}
\setmathfont{Latin Modern Math}
\usepackage{enumitem,needspace,xurl,hyperref}
\setlength{\parindent}{0pt}
\setlength{\parskip}{5pt}
\setlength{\emergencystretch}{3em}
\newcommand{\N}{\mathbb{N}}
\newcommand{\Z}{\mathbb{Z}}
\newcommand{\Q}{\mathbb{Q}}
\newcommand{\R}{\mathbb{R}}
\newcommand{\C}{\mathbb{C}}
\newcommand{\F}{\mathbb{F}}
\newcommand{\A}{\mathbb{A}}
\newcommand{\PP}{\mathbb{P}}
\newcommand{\E}{\mathbb{E}}
\newcommand{\eps}{\varepsilon}
\newcommand{\dd}{\,\mathrm{d}}
\begin{document}
\section*{UnsolvedMath ID: 30004708}
\textbf{Source problem code:} catalogue entry
\subsection{Definitions and mathematical statement}
Under the conventional interpretation of a hyperbolic metric as a complete Riemannian metric of constant sectional curvature \(-1\), fix an integer \(n\geq 2\) and a real number \(v>0\). Let
\[
\mathcal H_n(v)=\{(M,g): M\text{ is a connected closed hyperbolic }n\text{-manifold},\ \operatorname{Vol}_g(M)\leq v\}/\text{isometry}.
\]
Here closed means compact without boundary. No orientability restriction is stated in the source. The curvature normalization and dimension range above are conventional interpretations that require confirmation against the cited report.

For \((M,g)\in\mathcal H_n(v)\), define its systole by
\[
\operatorname{sys}(M,g)=\min\{\ell_g(\gamma):\gamma\text{ is a nonconstant closed geodesic in }M\}.
\]
Provisionally interpret its kissing number as
\[
K(M,g)=\#\{\gamma:\gamma\text{ is an unoriented primitive closed geodesic},\ \ell_g(\gamma)=\operatorname{sys}(M,g)\},
\]
where a primitive geodesic is not an iterate of a shorter closed geodesic, and geodesics differing only by parametrization or reversal are counted once. The source does not specify this counting convention, so this definition requires confirmation.

Let \(\Delta_g=-\operatorname{div}_g\nabla_g\) be the nonnegative Laplace--Beltrami operator on real-valued functions. Write its eigenvalues, counted with multiplicity, as
\[
0=\lambda_0(M,g)<\lambda_1(M,g)\leq\lambda_2(M,g)\leq\cdots,
\qquad
m_1(M,g)=\dim_{\R}\ker(\Delta_g-\lambda_1(M,g)I).
\]
The choice of the scalar Laplace--Beltrami operator is the conventional reading of the supplied wording.

The problem is to determine the extremal values of these four invariants over \(\mathcal H_n(v)\). More precisely, when this class is nonempty, the separate optimization interpretation asks for
\[
\sup_{(M,g)\in\mathcal H_n(v)}\operatorname{sys}(M,g),\qquad
\sup_{(M,g)\in\mathcal H_n(v)}K(M,g),\qquad
\sup_{(M,g)\in\mathcal H_n(v)}\lambda_1(M,g),\qquad
\sup_{(M,g)\in\mathcal H_n(v)}m_1(M,g),
\]
and for the manifolds attaining each supremum, when it is attained. No existence of maximizers is asserted. The source's phrase ``determine or maximize'' does not specify whether separate extrema, simultaneous optimization, or a broader classification is intended; the displayed formulation records the separate optimization interpretation.

For an orientable closed hyperbolic surface of genus \(g\geq2\),
\[
\operatorname{Area}(M)=4\pi(g-1).
\]
Thus prescribing its area is equivalent to prescribing its genus, whereas imposing an area upper bound only bounds the possible genera. For a nonorientable closed hyperbolic surface of nonorientable genus \(h\geq3\), the corresponding formula is \(\operatorname{Area}(M)=2\pi(h-2)\). Consequently, the source's assertion that volume is equivalent to fixing genus in dimension two needs both an orientability convention and a distinction between prescribed volume and a volume upper bound.
\subsection{Short English statement}
Determine the largest possible systole, number of shortest closed geodesics, first positive scalar Laplace eigenvalue, and its multiplicity among connected closed hyperbolic manifolds with volume at most \(v\). The draft interprets these as four separate optimization problems. Counting conventions, orientability, and the intended relation between bounded area and fixed genus require review.
\subsection{Sources}
\begin{itemize}
\item Supplied local UnsolvedMath registry record: \texttt{data/unsolvedmath/problems.json}, permanent ID \(30004708\), source code \texttt{catalogue entry}, title ``Extremal Invariants of Bounded-Volume Hyperbolic Manifolds,'' in the Oberwolfach Reports Open Problems collection. No individual proposer is supplied; the proposed-year field is \(2021\).
\item Supplied source citation: \emph{Enumerative Geometry of Surfaces (2022)}, Oberwolfach Reports, DOI \texttt{10.4171/owr/2021/28}, pp.\ 1463--1518. \url{https://doi.org/10.4171/owr/2021/28}. The background describes this collection as drawn from workshop reports and their open-problem discussions. No passage locating this particular problem within the report is supplied.
\item Supplied literature source labeled ``Fortier Bourque--Petri'': \url{https://arxiv.org/abs/2302.02540}. No title or specific supporting result is supplied.
\item Literature assessment dated \texttt{2026-08-22}: the supplied curation labels the problem open and states that general optimization of these invariants at fixed genus and volume remains unresolved, with bounds and isolated cases rather than a full classification. Its listed sources are the Oberwolfach report and Fortier Bourque--Petri links above. This is supplied dated curation, not an independently verified literature assessment; the record expressly requests independent verification of time-sensitive status and cited claims.
\item The supplied statement-status label is \texttt{corrected\_verified}; its accompanying explanation describes a publication-format audit repairing notation or typography while preserving mathematical content. This label does not establish mathematical verification of the present formulation.
\end{itemize}
\end{document}
